\section{The completed row energy}
\label{gex:row-energy}

The direct estimate uses the energy of the completed rows after
reflection. We recompute its exponent from \cite[Lemma~14.3]{OAI26};
the assembled estimate of \cite[Lemma~15.1]{OAI26} uses
\(\ell=1/6\).

Fix the target datum and slot system before \(Z\).  For a rescaled
subset \(J\) and tuple \((p_i)_{i\in J}\), write
\begin{equation}
\begin{gathered}
 d=\sum_{i\in J}\ell_i,\qquad \ell'=\ell-d,\qquad M'=1-\ell-2d,\\
 q_{p_J}=Z^dr_J,\qquad
 X'=Z^{l_x-d}/r_J,\qquad Y'=Z^{l_y-d}/r_J,\qquad
 Q=q_{b_*}X'Y'\asymp Z^{M'}.
\end{gathered}
\label{gex:rescaled-low-scales}
\end{equation}
Here \(r_J\) lies in a fixed positive compact interval.  We use exactly
the row \(B_{m,\sigma}^{J}\) defined in \cite[\S15.1]{OAI26}, with
ray expansion from \cite[\S6]{OAI26}.  Its surviving prime weights are
\(\prod_{i\notin J}\overline{\eta(p_i)}W_i(q_{p_i}/P_i)\), summed
inside the row before squaring.  Fixed window bounds enter the constant.

These product weights are independent of the row and dual variables.
The external squarefree label is one.  Additional punctures are fixed,
and the moving character \(\chi_\bullet(m)\) retains its zeros.
Local symbol powers are extended by zero on nonunits also at exponent
zero.  All nonresidual local primes are frozen; residual rows
and active marked products are separated into classes modulo the full
fixed reflection modulus.  We retain the original zero extensions and
dual support, including shared primes.  The row ball includes every
nonzero element, with its unit and bad-prime sectors.

\begin{lemma}
\label{gex:completed-low}
For a fixed disjoint slot system with \(0\le\ell\le1/5\), and every
\(\varepsilon>0\),
\begin{equation}
 \sum_{0<q_m\ll Q}|B_{m,\sigma}^{J}(Z)|^2
 \ll_{\varepsilon,\mathrm{fixed}}
 Z^{M'+(5\ell-1+d)_+/4+\varepsilon}.
\label{gex:completed-low-bound}
\end{equation}
The constant and threshold are uniform in the rescaled tuple.
They depend on the fixed datum and finitely many window seminorms.
\end{lemma}

\begin{proof}
Use the block variables of \cite[\S14.4]{OAI26}: \(O,H\) are the
powerful and residual row lengths; \(A_0\) is the active frozen
radical length; \(N_0,B_0,S_0\) describe the forced and small local
factors.  The active marked length is \(z_a\), and the remaining
dual lengths are \(v,\ell_b,e_\lambda\).  Put
\begin{equation}
 T_d=2H+2A_0+2z_a-N_*-N_0-3B_0,\qquad
 u=\min\{v,z_a,(v+z_a)/3\}.
\label{gex:retained-low-parameters}
\end{equation}
The imported bound for a retained block, after summing its frozen
parts, is \(Z^{E_{\rm ref}+\varepsilon_0}\), where
\begin{equation}
\begin{aligned}
 E_{\rm ref}={}&O/2+\max(H,v+\ell_b)-S_0-B_0+z_a-u-\ell_b
       -2e_\lambda/3\\
 &-\tfrac12(T_d-v-3\ell_b-e_\lambda)_+.
\end{aligned}
\label{gex:reflected-low-input}
\end{equation}
It is uniform on bounded length ranges, with finite seminorms and
polynomial norm-height cost.  The powerful-part count \(O/2\) is
already included.  Unretained whole blocks are arbitrarily power-small,
with their finite kernel order chosen before the polynomial row,
tuple and local multiplicities are summed.

First take \(N_*=1+\ell'\) and suppress the small annular errors.
The actual low source gives
\begin{equation}
\begin{gathered}
 \Delta_H=M'-O-H\ge0,\quad 2A_0\le O,\quad 0\le N_0\le A_0,\quad
 S_0,B_0\ge0,\\
 0\le z_a\le\ell',\quad v,\ell_b,e_\lambda\ge0,\quad
 v+3\ell_b+e_\lambda\le T_d.
\end{gathered}
\label{gex:actual-low-constraints}
\end{equation}
Since \(M'+\ell'-1=-3d\),
\begin{equation}
 T_d-(H-3d)
 =H-M'+2A_0+2(z_a-\ell')-N_0-3B_0\le0.
\label{gex:low-row-branch}
\end{equation}
Hence \(\max(H,v+\ell_b)=H\).

If \(u=z_a\), \eqref{gex:reflected-low-input} gives
\(E_{\rm ref}\le O/2+H\le M'-\Delta_H\).
Otherwise \(u\ge v/2\): when the minimum is \((v+z_a)/3\), its
comparison with \(z_a\) gives \(v\le2z_a\).
Writing \(y=v+3\ell_b+e_\lambda\), we obtain
\[
 u+\ell_b+2e_\lambda/3+\tfrac12(T_d-y)_+
 \ge y/4+\tfrac12(T_d-y)_+\ge T_d/4.
\]
Thus
\begin{align}
 E_{\rm ref}
 &\le O/2+H-S_0-B_0+z_a-T_d/4\notag\\
 &\le M'+(1+3\ell'-2M'-2\Delta_H)/4\notag\\
 &=M'+(5\ell-1+d-2\Delta_H)/4.
\label{gex:second-low-branch}
\end{align}
Here \(z_a\le\ell'\) and \(2A_0-N_0+B_0+4S_0\ge0\).
Both branches therefore give
\begin{equation}
 E_{\rm ref}\le M'+(5\ell-1+d)_+/4.
\label{gex:generalized-low-exponent}
\end{equation}

It remains to restore the Gaussian row.  For
\(\varrho_i=q_{p_i}/P_i\) and \(r=q_{cn^3}/Z^{1+\ell'}\), the
annulus \(r=e^kx\) has the joint profile
\[
 \chi(\log x)W_{\rm G}\!\left(
       e^kx/\prod_{i\notin J}\varrho_i\right)
 \prod_{i\notin J}W_i(\varrho_i),
 \qquad W_{\rm G}(y)=\frac{e^{-(\log y)^2/4}}{2\sqrt\pi}.
\]
Its homogeneous seminorms are summable against every fixed
\(e^{A|k|}\).  Lemmas~4.5--4.6 of \cite{OAI26} discard whole
annuli with \(|k|>\kappa_{\rm G}\log Z+C\) at arbitrary power
saving and separate each remaining profile through a Fourier
density common to all rows and tuples.  The prime weights remain
product weights, and \(q_c^{-1/2}q_n^{-1}\) is unchanged.
The Gaussian decay absorbs the polynomial counting factors before
these annular tails are discarded.

The completed length is now \(1+\ell'+k/\log Z\).
Dyad comparisons and the bounded negative ramified range add
\(O_{\rm fixed}(1/\log Z)\) errors.  Since the preceding expressions
are piecewise affine in finitely many lengths, their total error
is \(O(\kappa_{\rm G}+\tau_{\rm ref}+1/\log Z)\), where
\(\tau_{\rm ref}>0\) is the retention slack in Lemma~14.3.
Choose these allowances inside \(\varepsilon\), then enlarge the
threshold.  The kernel amplitude is extracted at the actual row
width before positive sieve enlargement.  Summing the logarithmic
dyads and finite sectors, followed by Minkowski through the common
density and Gaussian annuli, proves the lemma.
\end{proof}
